\chapter{Parallel and Multithreaded Algorithms}

Every algorithm you have studied so far has lived a lonely life. It does one thing
at a time. Step follows step follows step, a single thread of attention crawling
through the work like a single reader going through a book line by line. That
picture matched the machines of the last century, where one processor did one
operation per tick. But the machine on your desk right now is not like that. It
has many cores --- maybe four, maybe sixteen, maybe more --- and each of them can
be doing useful work at the very same instant. A phone has them. A laptop has
them. A data center has hundreds of thousands of them. If your algorithm insists
on doing one thing at a time, all but one of those cores sit idle, sipping
electricity and accomplishing nothing.

This chapter is about teaching algorithms to use many hands at once. You already
know the most important idea. Divide-and-conquer \emph{naturally} produces
independent subproblems --- ``sort the left half'' and ``sort the right half''
have nothing to say to each other until the very end --- and independent work is
exactly the kind of work that can be done in parallel. Much of what follows is
the divide-and-conquer you already know, now allowed to run its branches side by
side.

Parallelism comes with a new bookkeeping problem, though. When work happens at
the same time, ``how long does this take?'' splits into two different questions,
and we need two different cost measures to answer them. Learning to think with
both numbers at once is the real content of this chapter.

\begin{focusbox}
By the end of this chapter you should be able to: (1) describe an algorithm in
the \emph{fork-join} model using \textbf{spawn}, \textbf{sync}, and
\textbf{parallel for}; (2) compute its \emph{work} $T_1$ and its \emph{span}
$T_\infty$, and read off its \emph{parallelism} $T_1/T_\infty$; (3) bound how
fast it runs on $P$ processors using the work and span laws and Brent's theorem;
(4) parallelize loops, matrix multiplication, and merge sort, and analyze each by
solving a recurrence; and (5) avoid the classic traps --- counting work as if it
were time, and forgetting that a \emph{race} can silently corrupt your answer.
\end{focusbox}

\section{Two ways for work to happen at once}

Before any algorithm, fix the picture in your head. There are two flavors of
``doing things at the same time,'' and confusing them causes endless grief.

\begin{intuitionbox}
\textbf{Parallelism} is many workers doing the \emph{same job faster} by
splitting it up. Ten people painting one fence finish in a tenth of the time.
This chapter is about parallelism: we have a result to compute, and we want it
sooner by spreading the work across processors.

\textbf{Concurrency} is many independent activities that happen to be alive at
once --- a download, a clock, a spell-checker --- coordinated so they do not step
on each other. Concurrency is about \emph{managing} simultaneity; parallelism is
about \emph{exploiting} it for speed. They share machinery, but their goals are
different.
\end{intuitionbox}

We will also make a clean simplifying assumption, the one CLRS makes, that keeps
the analysis tractable: we imagine a friendly scheduler that we never have to
program. We describe \emph{what may run in parallel}, and an unseen runtime
system decides \emph{which core runs what, and when}. This separation lets us
reason about an algorithm's potential speed without drowning in the details of
any particular machine.

\begin{historybox}
This material is CLRS (4th edition) Chapter~26, ``Parallel Algorithms.'' The
fork-join / work-span model presented there grew out of the \emph{Cilk} project at
MIT in the 1990s (Blumofe, Leiserson, and others), whose runtime first made
\textbf{spawn} and \textbf{sync} both practical and provably efficient. The
work--span style of analysis traces back further, to the \emph{circuit} model of
parallel computation and to Richard Brent's 1974 theorem, which we meet at the end
of this chapter. If you continue to CLRS, you will find the same notation
($T_1$, $T_\infty$, parallelism) used there, so the vocabulary you build here
transfers directly.
\end{historybox}

\section{The fork-join model}

We need a way to \emph{write down} an algorithm that says ``these two pieces may
run at the same time.'' We use three keywords. They are deliberately few.

\begin{focusbox}
\textbf{spawn} --- placed before a procedure call, it says: ``you \emph{may} run
this call in parallel with the statements that come after it.'' The caller does
\emph{not} wait; it keeps going. (It is permission, not a command --- the
scheduler is free to run it serially if no core is free.)

\textbf{sync} --- ``wait here until every call I have spawned so far has
finished.'' This is the join. After a sync, every spawned result is available.

\textbf{parallel for} --- a loop whose iterations may all run at the same time
because they do not depend on one another.
\end{focusbox}

The name \emph{fork-join} is the shape these make. At a \textbf{spawn} the thread
of control \emph{forks} into two strands; at a \textbf{sync} the strands
\emph{join} back into one. Crucially there is an implicit sync at the end of every
procedure: a procedure never returns until all the work it spawned has completed.
That single rule keeps the model sane --- you can never accidentally use a result
before the worker that produced it is done.

Here is the canonical first example, computed two ways. Recall the Fibonacci
numbers $F_0=0$, $F_1=1$, $F_n=F_{n-1}+F_{n-2}$. The doubly-recursive program
below is a \emph{terrible} way to compute Fibonacci (it does exponential work),
but it is the perfect specimen for studying parallel structure, because its two
recursive calls are independent.

\begin{algorithm}[H]
\caption{\textsc{P-Fib}$(n)$}
\KwIn{A nonnegative integer $n$.}
\KwOut{The $n$th Fibonacci number $F_n$.}
\If{$n \le 1$}{\Return $n$}
$x \gets$ \textbf{spawn} \textsc{P-Fib}$(n-1)$\tcp*{may run in parallel}
$y \gets$ \textsc{P-Fib}$(n-2)$\tcp*{the caller does this meanwhile}
\textbf{sync}\tcp*{wait for the spawned call}
\Return $x + y$\;
\end{algorithm}

Read it slowly. We \textbf{spawn} the computation of $F_{n-1}$ and store its
eventual result in $x$. We do \emph{not} wait. The caller immediately marches on
and computes $F_{n-2}$ into $y$ itself. Only at the \textbf{sync} do we pause to
let the spawned strand catch up. After the sync, both $x$ and $y$ hold real
values, and we may add them. Drop the sync and we might read $x$ before it was
written --- a bug we name and dread shortly.

\begin{trapbox}
\textbf{spawn does not mean ``start a thread now and run it.''} It means ``it is
\emph{permitted} to run this in parallel.'' On a one-core machine, or when every
core is busy, the scheduler may simply run the spawned call serially, exactly as
an ordinary call. This is why a correct fork-join program must produce the same
answer whether or not anything actually runs in parallel. If your answer depends
on parallelism happening, you have written a bug, not an algorithm.
\end{trapbox}

\subsection{The computation as a graph}

The cleanest way to think about a fork-join computation is as a directed graph
that the execution traces out. Each vertex is a \emph{strand}: a maximal run of
instructions with no spawn, sync, or return in the middle. An edge $u \to v$ means
``$u$ must finish before $v$ may begin.'' A spawn creates a vertex with two
out-edges (the spawned strand and the continuation); a sync creates a vertex with
several in-edges that all must arrive before it fires. Because work cannot depend
on itself, this graph is a \emph{directed acyclic graph} (a DAG);
Figure~\ref{ch35:fork-join-dag} shows the smallest interesting one.

\begin{figure}[h]
\centering
\begin{tikzpicture}[
  >=Stealth, node distance=10mm,
  strand/.style={circle,draw=sourceblue,fill=blue!8,minimum size=6mm,inner sep=0pt,font=\small},
]
\node[strand] (a) {$a$};
\node[strand,below left=8mm and 12mm of a] (b) {$b$};
\node[strand,below right=8mm and 12mm of a] (c) {$c$};
\node[strand,below=20mm of a] (d) {$d$};
\draw[->,thick] (a) -- node[above left,font=\scriptsize] {spawn} (b);
\draw[->,thick] (a) -- node[above right,font=\scriptsize] {continue} (c);
\draw[->,thick] (b) -- (d);
\draw[->,thick] (c) -- node[right,font=\scriptsize] {sync} (d);
\node[align=center,font=\scriptsize,below=2mm of d] {join};
\end{tikzpicture}
\caption{One spawn-and-sync as a DAG. Strand $a$ forks into the spawned strand
$b$ and the continuation $c$; both must finish before the join strand $d$ may
run. The graph encodes \emph{dependencies}, not a schedule.}
\label{ch35:fork-join-dag}
\end{figure}

This graph is the object we measure. Everything in the analysis below is a
statement about it, and two numbers describe it completely for our purposes.

\section{Work and span: the two numbers that matter}

We now reach the heart of the chapter. There are exactly two questions you can
ask about a parallel computation that do not depend on the messy details of a
real machine, and each gives a number.

\begin{intuitionbox}
\textbf{Work} $T_1$: how much \emph{total} computation is there? Equivalently,
how long would it take on \emph{one} processor, which must do every strand
one after another. It is the sum of the running times of all the strands --- the
total area of the job, ignoring that some of it could overlap.

\textbf{Span} $T_\infty$: how long is the \emph{longest chain of dependencies}?
Equivalently, how long would it take on \emph{infinitely many} processors, where
the only thing that can ever slow you down is waiting for something that genuinely
had to happen first. It is the length of the longest path through the DAG --- the
``critical path.''
\end{intuitionbox}

The subscripts are a memory aid: $T_1$ is the time on $1$ processor, and
$T_\infty$ is the time on $\infty$ processors. The work is the \emph{area}; the
span is the \emph{depth}. Hold both pictures: a wide, shallow computation has lots
of work but a short span (very parallelizable); a tall, skinny one has a long span
no matter how many processors you throw at it.

\begin{walkbox}
Picture a relay race versus a choir. In a relay, runner two cannot start until
runner one finishes: the work (total distance) equals the span (longest chain), so
extra runners standing around do not help --- the dependencies are a single line.
In a choir, a hundred singers all sing at once: the work is a hundred voice-parts,
but the span is just the length of \emph{one} song. As you walk, classify
everyday tasks as ``relay-shaped'' or ``choir-shaped.'' Which kind benefits from
more hands? That intuition \emph{is} the work--span distinction.
\end{walkbox}

\subsection{The two laws}

These two numbers immediately bound the running time on a real machine with $P$
processors, which we write $T_P$. Two laws do the work, and both are obvious once
stated.

\begin{formalbox}
\textbf{The work law.} $P$ processors can do at most $P$ units of work per time
step, so to finish $T_1$ work takes at least $T_1/P$ steps:
$$T_P \ge \frac{T_1}{P}.$$
\textbf{The span law.} Adding processors can never make you faster than the
critical path, which must be walked one strand after another no matter what:
$$T_P \ge T_\infty.$$
\end{formalbox}

Read the work law as ``you cannot finish faster than sharing the total work
perfectly,'' and the span law as ``you cannot finish faster than the longest
unavoidable wait.'' The actual time $T_P$ is squeezed from below by both. A good
parallel algorithm is one where these two lower bounds are also roughly
achievable --- which is exactly what Brent's theorem, coming up, will promise.

\subsection{Speedup and parallelism}

With both laws in hand, we can define the quantities we have been circling.

\begin{definition}[Speedup and parallelism]\label{ch35:def-speedup}
The \emph{speedup} of a computation on $P$ processors is the ratio
$T_1 / T_P$ --- how many times faster the $P$-processor run is than the serial
run. The \emph{parallelism} of the computation is
$$\frac{T_1}{T_\infty},$$
the work divided by the span. It is a single number attached to the algorithm
itself, with no $P$ in it.
\end{definition}

By the work law, $T_1/T_P \le P$: with $P$ processors you can never go more than
$P$ times faster. When the speedup equals $P$ we call it \emph{perfect} (or
\emph{linear}) speedup. By the span law, $T_1/T_P \le T_1/T_\infty$: the speedup
can never exceed the parallelism, no matter how many processors you buy. So
parallelism is a ceiling. If an algorithm has parallelism $1000$, then putting
more than about $1000$ processors on it is pouring water into a full glass --- the
extra cores cannot help, because there is not enough independent work to feed
them. We say a computation has \emph{slack} when $P \ll T_1/T_\infty$, meaning
there is comfortably more parallelism available than processors to use it, and in
that happy regime good schedulers achieve near-perfect speedup.

\begin{focusbox}
Three slogans to carry forever:
\begin{itemize}
\item \textbf{Work} is total area; you minimize it to avoid \emph{wasting}
computation.
\item \textbf{Span} is critical-path depth; you minimize it to avoid
\emph{serializing} computation.
\item \textbf{Parallelism} $=$ work $/$ span is the most processors you can ever
usefully employ.
\end{itemize}
The art is keeping work low (so you are not slower than the serial program) while
\emph{also} keeping span low (so the parallelism is high). These two goals can
fight, and resolving the fight is what algorithm design in this chapter is about.
\end{focusbox}

\subsection{Worked example: the work and span of \textsc{P-Fib}}

Let us compute the two numbers for \textsc{P-Fib}$(n)$. Spawn, sync, and the
bookkeeping cost only constant time, so the recurrences look just like the serial
ones --- except that \emph{span} gets to take the longer of the two recursive
branches instead of their sum, because those branches run in parallel.

\textbf{Work} $T_1(n)$. On one processor, spawn is just an ordinary call, so the
work is exactly the serial running time of double-recursive Fibonacci:
$$T_1(n) = T_1(n-1) + T_1(n-2) + \Th(1).$$
This is the Fibonacci recurrence itself (plus constants), so $T_1(n) =
\Th(F_n) = \Th(\phi^n)$ where $\phi = (1+\sqrt5)/2 \approx 1.618$. The work is
exponential. (Again: this is a bad way to compute Fibonacci; we are using it only
to study structure.)

\textbf{Span} $T_\infty(n)$. The two recursive calls run \emph{in parallel}, so
the critical path goes through whichever one is longer, not through their sum.
Since $T_\infty(n-1) \ge T_\infty(n-2)$, the longer branch is the $(n-1)$ one:
$$T_\infty(n) = \max\big(T_\infty(n-1),\, T_\infty(n-2)\big) + \Th(1)
= T_\infty(n-1) + \Th(1) = \Th(n).$$
So the span is merely \emph{linear}. The critical path is the chain $n \to n-1 \to
n-2 \to \cdots \to 1$, one step per level.

\textbf{Parallelism.}
$$\frac{T_1(n)}{T_\infty(n)} = \frac{\Th(\phi^n)}{\Th(n)} = \Th\!\left(\frac{\phi^n}{n}\right).$$
This grows enormously with $n$. For $n=30$ it is already in the millions. There is
a vast amount of independent work here --- which is the silver lining of all that
wasteful recomputation. The lesson is sharp: high parallelism and low total work
are \emph{different} virtues, and an algorithm can have one without the other.

\begin{labbox}
Trace \textsc{P-Fib}$(4)$ by hand and build its DAG. The call $P(4)$ spawns
$P(3)$ and runs $P(2)$; $P(3)$ spawns $P(2)$ and runs $P(1)$; each $P(2)$ spawns
$P(1)$ and runs $P(0)$. Count the calls (that is your work, up to constants):
the number of calls is $2F_{n+1}-1$, which for $n=4$ is $2\cdot 5 - 1 = 9$ calls.
Now find the longest dependency chain from the top call down to a leaf and back up
through the adds: it has length proportional to $n=4$, not to $9$. Seeing ``$9$
calls but a depth-$4$ critical path'' in front of you (Figure~\ref{ch35:pfib-tree})
makes the work--span gap concrete in a way no formula can.
\end{labbox}

\begin{figure}[h]
\centering
\begin{tikzpicture}[
  >=Stealth, level distance=12mm, sibling distance=17mm,
  every node/.style={circle,draw=sourceblue,fill=blue!8,inner sep=1pt,
    minimum size=7mm,font=\scriptsize},
  edge from parent/.style={draw,->,gray}
]
\node {$P4$}
  child {node {$P3$}
    child {node {$P2$}
      child {node {$P1$}}
      child {node {$P0$}}}
    child {node {$P1$}}}
  child {node {$P2$}
    child {node {$P1$}}
    child {node {$P0$}}};
\end{tikzpicture}
\caption{The call tree of \textsc{P-Fib}$(4)$: $9$ calls in all (that is the
work, up to constants), but the longest root-to-leaf chain
$P4 \to P3 \to P2 \to P1$ has only $4$ nodes. The critical path runs down this
spine of left-most (largest-argument) children and back up through the additions,
so the span is $\Th(n)$ while the work is $\Th(\phi^n)$.}
\label{ch35:pfib-tree}
\end{figure}

\section{Scheduling: turning the two numbers into a real running time}

We have lower bounds ($T_P \ge T_1/P$ and $T_P \ge T_\infty$). Can we
\emph{achieve} something close to them on a real machine, with a scheduler that
does not know the future? We can, and a simple greedy rule is enough.

\begin{definition}[Greedy scheduler]\label{ch35:def-greedy}
A \emph{greedy} scheduler, at every time step, assigns as many ready strands to
processors as it can. A step is \emph{complete} if all $P$ processors are kept
busy; it is \emph{incomplete} if fewer than $P$ ready strands exist and some
processors idle. Greedy simply never leaves a processor idle when there is ready
work for it.
\end{definition}

\begin{theorem}[Greedy scheduling bound]\label{ch35:thm-greedy}
On $P$ processors, a greedy scheduler runs any computation with work $T_1$ and
span $T_\infty$ in time
$$T_P \le \frac{T_1}{P} + T_\infty.$$
\end{theorem}

\begin{intuitionbox}
Here is the idea. Classify each time step as \emph{complete} (all $P$ processors
busy) or \emph{incomplete} (some idle). There cannot be many complete steps,
because each one chews through $P$ units of work and there are only $T_1$ units
total --- so at most $T_1/P$ complete steps. There cannot be many incomplete
steps either: every incomplete step must make progress along the critical path
(it ran every ready strand, so the next strand on the longest remaining path
became ready), and the critical path has only $T_\infty$ strands --- so at most
$T_\infty$ incomplete steps. Add the two buckets: total steps
$\le T_1/P + T_\infty$.
\end{intuitionbox}

\begin{formalbox}
\textbf{Proof sketch.} Count complete steps. Each uses $P$ work, and total work is
$T_1$, so there are at most $\floor{T_1/P}$ of them. Count incomplete steps.
Consider the subgraph of strands not yet executed; its longest path starts with a
ready strand. On an incomplete step the scheduler runs \emph{all} ready strands
(there were fewer than $P$), so in particular it runs the first strand of that
longest path, shortening the remaining critical path by one. The remaining
critical path starts at length $T_\infty$ and cannot drop below $0$, so there are
at most $T_\infty$ incomplete steps. Every step is one or the other, giving
$T_P \le T_1/P + T_\infty$. $\qquad\blacksquare$
\end{formalbox}

This bound, in this clean form, is essentially Graham's and Brent's, and it is the
workhorse of the whole chapter: \emph{greedy gets you the work term plus the span
term, and never worse.} Compare it to the lower bound
$T_P \ge \max(T_1/P,\, T_\infty)$. The upper and lower bounds differ by at most a
factor of $2$, so greedy is always within a factor of $2$ of optimal. We can do
even better when there is slack.

\begin{corollary}[Near-perfect speedup with slack]\label{ch35:cor-slack}
If $P \ll T_1/T_\infty$ (the parallelism comfortably exceeds the processor count),
then $T_\infty \ll T_1/P$, so the span term in
Theorem~\ref{ch35:thm-greedy} is negligible and
$$T_P \approx \frac{T_1}{P}, \qquad \text{speedup } \frac{T_1}{T_P} \approx P.$$
That is, when there is plenty of slack, a greedy scheduler delivers near-perfect
linear speedup.
\end{corollary}

\begin{aestheticbox}
This is why the work--span model is the model real systems use. A programmer
writing in Cilk, Intel TBB, Java's fork-join framework, OpenMP tasks, or Rust's
Rayon does \emph{not} schedule threads by hand. They express the spawn/sync
structure and let a provably-good greedy (work-stealing) scheduler do the rest.
The two numbers you compute on paper --- work and span --- are exactly the two
numbers that predict how the program will scale when the hardware team doubles the
core count next year. You are not doing abstract math; you are computing the
slope of a performance graph.
\end{aestheticbox}

\section{Analyzing parallel loops}

Most real parallel code is not double recursion; it is loops. A \textbf{parallel
for} declares that all iterations are independent and may run at once. But ``run
all $n$ at once'' is a lie on real hardware, and even in the model it hides a cost
we must account for. The runtime does not literally fork $n$ strands in one step;
it forks them with a \emph{recursive binary splitting} of the index range, like a
balanced tournament bracket. That hidden tree is what the analysis turns on.

Consider the simplest possible parallel loop: add two vectors,
$C[i] \gets A[i] + B[i]$ for $i = 1 \ldots n$.

\begin{algorithm}[H]
\caption{\textsc{Vector-Add}$(A, B, C, n)$}
\KwIn{Arrays $A$, $B$ of length $n$ and an output array $C$.}
\KwOut{$C[i] = A[i]+B[i]$ for all $i$.}
\textbf{parallel for} $i \gets 1$ \textbf{to} $n$\;
\Indp $C[i] \gets A[i] + B[i]$\;
\Indm
\end{algorithm}

\textbf{Work.} There are $n$ iterations, each doing $\Th(1)$ work, plus the
overhead of the splitting tree. The splitting tree that forks $n$ iterations has
$\Th(n)$ leaves and hence $\Th(n)$ internal nodes, so the bookkeeping is also
$\Th(n)$. Total work $T_1(n) = \Th(n)$. Good --- the same as the serial loop.

\textbf{Span.} This is where the subtlety lives. If the loop body has span
$\Th(1)$, the span of the whole loop is \emph{not} $\Th(1)$. To reach $n$
independent iterations the
runtime must spawn them through a balanced binary tree of depth $\Th(\log n)$
(Figure~\ref{ch35:loop-tree}). You
walk \emph{down} the tree (forking) and back \emph{up} (joining), so
$$T_\infty(n) = \Th(\log n) + \Th(1) = \Th(\log n).$$
The $\log n$ is the unavoidable price of distributing the work; the trailing
$\Th(1)$ is the body. If the body itself had span $s$, the loop's span would be
$\Th(\log n) + s$.

\begin{figure}[h]
\centering
\begin{tikzpicture}[
  >=Stealth, level distance=11mm, sibling distance=20mm,
  every node/.style={draw=quotientorange,fill=orange!8,circle,inner sep=1pt,minimum size=5mm,font=\scriptsize},
  edge from parent/.style={draw,->,thick}
]
\node {1--8}
  child {node {1--4}
    child {node {1--2}}
    child {node {3--4}}}
  child {node {5--8}
    child {node {5--6}}
    child {node {7--8}}};
\end{tikzpicture}
\caption{How a \textbf{parallel for} over $8$ iterations is really executed: a
balanced splitting tree (shown two levels deep) of depth $\Th(\log n)$. The leaves
are the iterations. Walking down and back up this tree is why the span carries a
$\log n$, even when each iteration is $\Th(1)$.}
\label{ch35:loop-tree}
\end{figure}

\textbf{Parallelism.} $T_1/T_\infty = \Th(n)/\Th(\log n) = \Th(n/\log n)$. Huge.
For a million-element vector that is on the order of $50{,}000$ --- you would need
fifty thousand processors before the parallelism ran dry. Parallel loops over big
arrays are about as good as parallelism gets.

\begin{trapbox}
Do not write $T_\infty = \Th(1)$ for a parallel loop just because each iteration
is constant time. The \emph{logarithmic} cost of forking the iterations is real
and shows up in the span every single time. Forgetting it is the most common
mistake in span analysis. The rule of thumb: \emph{a parallel for adds $\log
(\text{number of iterations})$ to the span.}
\end{trapbox}

\begin{walkbox}
Why a \emph{tree} and not a single fork into $n$? Imagine one strand trying to
spawn a million iterations by itself, one spawn at a time. That spawning loop is
serial --- it takes a million steps just to hand out the work, so the span would be
$\Th(n)$ and the parallelism would collapse to $\Th(1)$. The balanced tree lets
the spawning itself be parallel: each node spawns only two children, so after
$\log n$ levels all $n$ leaves exist. On your walk, convince yourself that
\emph{distributing} work is itself work that must be parallelized, or it becomes
the bottleneck.
\end{walkbox}

\section{A new hazard: races}

Sequential code cannot have a race. Parallel code can, and a race is the most
insidious bug in this chapter because the program usually \emph{works} --- until,
one run in a thousand, it silently gives a wrong answer.

\begin{definition}[Determinacy race]\label{ch35:def-race}
A \emph{determinacy race} occurs when two strands that are \emph{not} ordered by
the dependency DAG both access the same memory location and at least one of them
writes to it. Because the scheduler may interleave the two strands in either
order, the result can differ from run to run.
\end{definition}

\begin{trapbox}
Here is the classic poison. Suppose you tried to sum an array in parallel by
having every iteration add into a shared variable:
\begin{center}
\textbf{parallel for} $i \gets 1$ \textbf{to} $n$: \quad $total \gets total + A[i]$
\end{center}
The statement $total \gets total + A[i]$ is really three steps: \emph{read} $total$,
\emph{add} $A[i]$, \emph{write} $total$. If two iterations read the old value
\emph{before} either writes, one of the additions is lost. Run it on $A=[1,1,1,1]$
and you might get $4$, or $3$, or $2$, depending on timing. The two writing
strands are unordered in the DAG, both touch $total$, and at least one writes ---
the definition of a race. The fix is to make the updates not collide: split the
array, sum the halves into \emph{separate} variables in parallel, sync, then add
the two partial sums. That is a parallel \emph{reduction}
(Figure~\ref{ch35:reduction-tree}), and it has span $\Th(\log n)$ with no race.
\end{trapbox}

\begin{figure}[h]
\centering
\begin{tikzpicture}[>=Stealth, font=\scriptsize,
  every node/.style={draw=targetgreen,fill=green!8,rounded corners,
    minimum width=8mm,minimum height=5mm,inner sep=2pt}]
% leaves
\node (a) at (0,0) {$A_1$};
\node (b) at (1.5,0) {$A_2$};
\node (c) at (3,0) {$A_3$};
\node (d) at (4.5,0) {$A_4$};
% level 1
\node (ab) at (0.75,1.1) {$+$};
\node (cd) at (3.75,1.1) {$+$};
% root
\node (abcd) at (2.25,2.2) {$+$};
\draw[->] (a) -- (ab); \draw[->] (b) -- (ab);
\draw[->] (c) -- (cd); \draw[->] (d) -- (cd);
\draw[->] (ab) -- (abcd); \draw[->] (cd) -- (abcd);
\node[draw=none,fill=none,right=2mm of abcd] {total};
\end{tikzpicture}
\caption{Race-free parallel summation (a \emph{reduction}). Pairs of partial sums
combine up a balanced tree; siblings are computed in parallel and each ``$+$''
node writes its own cell, so no two parallel strands ever touch the same location.
The tree has depth $\Th(\log n)$, giving span $\Th(\log n)$ with work $\Th(n)$.}
\label{ch35:reduction-tree}
\end{figure}

The model gives us a clean way to stay safe: \emph{a fork-join program is
race-free if every pair of strands that touch the same location, one of them
writing, is ordered by a spawn/sync dependency.} Put differently, never let two
\emph{parallel} strands write the same cell, or write-and-read the same cell. In
all the algorithms below we are careful to give each parallel branch its
\emph{own} disjoint piece of memory to write into --- the left half writes the
left half of the output, the right half writes the right half --- so no two
parallel strands ever collide. That discipline is what makes the model's clean
analysis legitimate.

\section{Parallel matrix multiplication}

Now we apply everything to a real, useful problem. Multiplying two $n \times n$
matrices, $C = A \cdot B$, sits at the center of scientific computing, graphics,
and machine learning, and it parallelizes beautifully because \emph{every output
entry can be computed independently}.

\subsection{The straightforward version}

Recall the definition $C[i][j] = \sum_{k=1}^{n} A[i][k]\,B[k][j]$. The $n^2$
output entries do not depend on each other, so compute them all in parallel.

\begin{algorithm}[H]
\caption{\textsc{P-Matrix-Multiply}$(A, B, C, n)$}
\KwIn{$n \times n$ matrices $A$, $B$; output matrix $C$.}
\KwOut{$C = A \cdot B$.}
\textbf{parallel for} $i \gets 1$ \textbf{to} $n$\;
\Indp \textbf{parallel for} $j \gets 1$ \textbf{to} $n$\;
\Indp $C[i][j] \gets 0$\;
\For{$k \gets 1$ \textbf{to} $n$}{$C[i][j] \gets C[i][j] + A[i][k]\cdot B[k][j]$}
\Indm\Indm
\end{algorithm}

Notice the inner $k$-loop is an \emph{ordinary} (serial) for, not a parallel one.
That is deliberate: it accumulates into the single cell $C[i][j]$, so making it
parallel would create exactly the summation race from the last section. Each
$(i,j)$ pair owns its own output cell, so the two outer parallel loops are
perfectly safe.

\textbf{Work.} Every cell does $\Th(n)$ multiply-adds, and there are $n^2$ cells,
so $T_1(n) = \Th(n^3)$ --- identical to the serial algorithm. We waste nothing.

\textbf{Span.} The outer parallel-for contributes $\Th(\log n)$; the nested
parallel-for inside it contributes another $\Th(\log n)$; and the serial inner
loop is a chain of length $\Th(n)$. The span is the sum of these along the
critical path:
$$T_\infty(n) = \Th(\log n) + \Th(\log n) + \Th(n) = \Th(n).$$
The serial $k$-loop dominates --- its length-$n$ chain is the critical path.

\textbf{Parallelism.} $T_1/T_\infty = \Th(n^3)/\Th(n) = \Th(n^2)$. For
$n = 1000$ that is a parallelism of a million. This simple algorithm already
scales superbly. Its only blemish is the $\Th(n)$ span coming from the serial dot
product, and we can shorten that by parallelizing the summation too, with a
divide-and-conquer multiply.

\subsection{Divide-and-conquer multiply, for a shorter span}

Split each $n \times n$ matrix into four $n/2 \times n/2$ blocks:
$$A = \begin{pmatrix} A_{11} & A_{12} \\ A_{21} & A_{22}\end{pmatrix},\quad
B = \begin{pmatrix} B_{11} & B_{12} \\ B_{21} & B_{22}\end{pmatrix},\quad
C = \begin{pmatrix} C_{11} & C_{12} \\ C_{21} & C_{22}\end{pmatrix}.$$
Block matrix multiplication gives eight smaller products:
$$\begin{aligned}
C_{11} &= A_{11}B_{11} + A_{12}B_{21}, & C_{12} &= A_{11}B_{12} + A_{12}B_{22},\\
C_{21} &= A_{21}B_{11} + A_{22}B_{21}, & C_{22} &= A_{21}B_{12} + A_{22}B_{22}.
\end{aligned}$$
The eight recursive multiplies are independent (each writes into a scratch area),
so spawn all eight; sync; then add the pairs in parallel.

\textbf{Work.} Eight recursive multiplies of size $n/2$, plus $\Th(n^2)$ for the
four matrix additions:
$$T_1(n) = 8\,T_1(n/2) + \Th(n^2).$$
By the master theorem, $a=8$, $b=2$, $\log_b a = 3$, and $n^{\log_b a} = n^3$
dominates $n^2$, so $T_1(n) = \Th(n^3)$ --- still cubic, as expected.

\textbf{Span.} The eight multiplies run in parallel, so the span takes \emph{one}
recursive call's span, not eight. The block additions can be done with parallel
fors over $n^2$ entries, contributing $\Th(\log n)$:
$$T_\infty(n) = T_\infty(n/2) + \Th(\log n).$$
Unrolling: $\Th(\log n) + \Th(\log(n/2)) + \cdots$ over $\log n$ levels, each
level's term at most $\Th(\log n)$, giving $T_\infty(n) = \Th(\log^2 n)$.

\textbf{Parallelism.} $T_1/T_\infty = \Th(n^3)/\Th(\log^2 n) = \Th(n^3/\log^2 n)$.
This is enormously larger than the $\Th(n^2)$ of the simple version --- the span
dropped from $\Th(n)$ all the way to $\Th(\log^2 n)$ by refusing to do that one
serial summation. Same work, dramatically shorter span: a textbook illustration of
the design goal.

\begin{complexitybox}
\textbf{Parallel matrix multiply, $n \times n$.}
\begin{center}
\begin{tabular}{lccc}
\toprule
Method & Work $T_1$ & Span $T_\infty$ & Parallelism \\
\midrule
Triple-loop (serial $k$) & $\Th(n^3)$ & $\Th(n)$ & $\Th(n^2)$ \\
Divide-and-conquer (8-way) & $\Th(n^3)$ & $\Th(\log^2 n)$ & $\Th(n^3/\log^2 n)$ \\
\bottomrule
\end{tabular}
\end{center}
Both do the same total work. The divide-and-conquer version trades nothing in work
for a vastly shorter critical path, hence far more parallelism.
\end{complexitybox}

\section{Parallel merge sort}

Sorting is the chapter's most satisfying case study, because the obvious
parallelization is good but not great, and fixing it teaches the deepest lesson in
the chapter: \emph{sometimes you must parallelize the combine step too.}

\subsection{The naive parallelization, and its disappointing span}

Merge sort divides, sorts the two halves, and merges. The two halves are
independent --- spawn them.

\begin{algorithm}[H]
\caption{\textsc{P-Merge-Sort}$(A, lo, hi)$}
\KwIn{Array $A$ and bounds $lo \le hi$.}
\KwOut{$A[lo \ldots hi]$ sorted in place.}
\If{$lo \ge hi$}{\Return}
$mid \gets \floor{(lo+hi)/2}$\;
\textbf{spawn} \textsc{P-Merge-Sort}$(A, lo, mid)$\tcp*{sort left half in parallel}
\textsc{P-Merge-Sort}$(A, mid+1, hi)$\tcp*{sort right half here}
\textbf{sync}\;
\textsc{Merge}$(A, lo, mid, hi)$\tcp*{serial merge: the bottleneck}
\end{algorithm}

\textbf{Work.} Two recursive sorts of size $n/2$ plus a serial merge of $\Th(n)$:
$$T_1(n) = 2\,T_1(n/2) + \Th(n) = \Th(n \log n),$$
by the master theorem ($a=b=2$, $\log_b a = 1$, the $n^1$ and the merge tie, so a
log factor appears). Same as serial merge sort --- no work wasted.

\textbf{Span.} The two recursive sorts run in parallel, so the span uses one of
them; but the \textsc{Merge} is serial, contributing $\Th(n)$:
$$T_\infty(n) = T_\infty(n/2) + \Th(n) = \Th(n).$$
That last recurrence solves to $\Th(n)$ because the $\Th(n)$ merge at the
\emph{top} level dominates: $n + n/2 + n/4 + \cdots = \Th(n)$.

\textbf{Parallelism.} $T_1/T_\infty = \Th(n\log n)/\Th(n) = \Th(\log n)$. That is
\emph{terrible}. A parallelism of only $\log n$ means even for a billion elements
you can usefully employ only about $30$ processors. The villain is plain: the
serial \textsc{Merge}, with its $\Th(n)$ span, single-handedly drags the whole
algorithm down. We sorted in parallel but then merged in a long serial line. To
get real parallelism, we must learn to \emph{merge in parallel}.

\subsection{Merging two sorted arrays in parallel}

How do you merge $X[1 \ldots n_1]$ and $Y[1 \ldots n_2]$ (both sorted) using many
hands? The serial merge walks a single pointer down the output, which is
inherently a chain of length $n$. We need a divide-and-conquer merge instead.

\begin{intuitionbox}
\textbf{The pivot trick.} Take the \emph{median} element of the larger array, say
$X[q]$ where $q$ is the middle index of $X$. Everything in the merged output that
is less than $X[q]$ comes from the front of $X$ \emph{and} the front of $Y$; we
find where $X[q]$ belongs in $Y$ by \emph{binary search}, splitting $Y$ at some
position $q'$. Now we have two completely independent merge problems --- merge
$X$'s left part with $Y$'s left part, and $X$'s right part with $Y$'s right part
--- and we know exactly where each result goes in the output. Spawn both. No race:
the left merge writes the left output, the right merge writes the right output.
Figure~\ref{ch35:pmerge-split} shows the split.
\end{intuitionbox}

\begin{figure}[h]
\centering
\begin{tikzpicture}[>=Stealth, font=\scriptsize]
% X array
\draw (0,1.4) rectangle (5,2.0);
\draw (2.4,1.4) rectangle (2.6,2.0);
\fill[sourceblue!25] (2.4,1.4) rectangle (2.6,2.0);
\node at (1.2,1.7) {$X$ left ($<X[q]$)};
\node at (3.8,1.7) {$X$ right};
\node[sourceblue] at (2.5,2.35) {$X[q]$};
\draw[sourceblue,->] (2.5,2.25) -- (2.5,2.02);
\node[left] at (-0.1,1.7) {$X$:};
% Y array
\draw (0,0) rectangle (5,0.6);
\draw[targetgreen,very thick] (3.1,0) -- (3.1,0.6);
\node at (1.5,0.3) {$Y$ left ($<X[q]$)};
\node at (4.0,0.3) {$Y$ right};
\node[targetgreen] at (3.1,-0.35) {split at $q'$};
\node[left] at (-0.1,0.3) {$Y$:};
% braces for left/right merge problems
\draw[decorate,decoration={brace,amplitude=4pt},gray]
  (-0.2,2.15) -- (-0.2,-0.15);
\node[gray,align=center] at (6.6,1.0)
  {left merge $\to Z$\\(both fronts)};
\node[gray,align=center] at (6.6,2.0) {};
\draw[gray,->] (5.05,1.0) -- (5.9,1.0);
\end{tikzpicture}
\caption{The \textsc{P-Merge} pivot trick. The median $X[q]$ of the longer array
$X$ is located inside $Y$ by binary search (position $q'$). Everything to the left
of both splits (less than $X[q]$) becomes one independent merge subproblem; the two
right parts become the other. The two subproblems write into disjoint output
regions, so they may be spawned without a race.}
\label{ch35:pmerge-split}
\end{figure}

\begin{algorithm}[H]
\caption{\textsc{P-Merge}$(X, p_1, r_1, Y, p_2, r_2, Z, p_3)$}
\KwIn{Sorted ranges $X[p_1 \ldots r_1]$, $Y[p_2 \ldots r_2]$; output array $Z$
starting at index $p_3$.}
\KwOut{The merged sorted sequence written into $Z[p_3 \ldots]$.}
$n_1 \gets r_1 - p_1 + 1$; \quad $n_2 \gets r_2 - p_2 + 1$\;
\If{$n_1 < n_2$}{swap the two input ranges so that $X$ is the longer one\;}
\If{$n_1 = 0$}{\Return\tcp*{both ranges empty}}
$q \gets \floor{(p_1 + r_1)/2}$\tcp*{median of the longer range $X$}
$q' \gets$ \textsc{Binary-Search}$(X[q], Y, p_2, r_2)$\tcp*{first index in $[p_2,r_2+1]$ with $Y[q'] \ge X[q]$}
$q_3 \gets p_3 + (q - p_1) + (q' - p_2)$\tcp*{output position for $X[q]$}
$Z[q_3] \gets X[q]$\;
\textbf{spawn} \textsc{P-Merge}$(X, p_1, q-1, Y, p_2, q'-1, Z, p_3)$\tcp*{left halves}
\textsc{P-Merge}$(X, q+1, r_1, Y, q', r_2, Z, q_3+1)$\tcp*{right halves}
\textbf{sync}\;
\end{algorithm}

Here \textsc{Binary-Search}$(x, Y, p_2, r_2)$ returns the first index $q'$ in the
range $[p_2,\, r_2+1]$ with $Y[q'] \ge x$; if no such element exists it returns
$r_2+1$. The elements $Y[p_2 \ldots q'-1]$ are exactly those strictly below $X[q]$,
so $q'-p_2$ of $Y$'s elements precede $X[q]$ in the output. Adding the $q-p_1$
elements of $X$ that precede $X[q]$ gives its output slot $q_3$. The left
recursive call then merges everything below $X[q]$ (the front of both ranges), the
right call merges everything at or above it, and the two write into disjoint
output regions $Z[p_3 \ldots q_3-1]$ and $Z[q_3+1 \ldots]$ --- so they cannot
race.

\textbf{Span of \textsc{P-Merge}.} Each call does a binary search ($\Th(\log n)$)
and then recurses on two independent subproblems in parallel. By always pivoting
on the median of the \emph{larger} array, we guarantee each recursive call handles
at most $3/4$ of the total elements. To see why, let $X$ be the longer range, so
$\abs{X} \ge n/2$ and $\abs{Y} \le n/2$. Splitting $X$ at its median sends at most
$\abs{X}/2$ of its elements to either side, and the worst case dumps \emph{all} of
$Y$ onto one side. So a recursive call gets at most
$\abs{X}/2 + \abs{Y} = (n-\abs{Y})/2 + \abs{Y} = n/2 + \abs{Y}/2 \le n/2 + n/4 =
3n/4$. So
$$T_\infty^{\text{merge}}(n) = T_\infty^{\text{merge}}(3n/4) + \Th(\log n)
= \Th(\log^2 n).$$
The recursion has depth $\log_{4/3} n = \Th(\log n)$, and each level costs
$\Th(\log n)$ for its binary search, multiplying to $\Th(\log^2 n)$.

\textbf{Work of \textsc{P-Merge}.} One can show (it takes a little care because
the split is uneven) that the work is still $T_1^{\text{merge}}(n) = \Th(n)$ ---
the binary searches sum to only $\Th(n)$ across the whole recursion, the same as
the serial merge up to constants.

\subsection{Putting it together: a genuinely parallel merge sort}

Now replace the serial \textsc{Merge} in \textsc{P-Merge-Sort} with
\textsc{P-Merge}. (A clean implementation sorts into an auxiliary array so the
parallel merge has separate input and output regions and stays race-free.) The
recurrences become:

\textbf{Work.} Same shape, $T_1(n) = 2\,T_1(n/2) + \Th(n) = \Th(n\log n)$. The
parallel merge does the same total work as the serial one, so the work is
unchanged. Good --- we are not paying for the speedup in wasted computation.

\textbf{Span.} Now the merge contributes only $\Th(\log^2 n)$, not $\Th(n)$:
$$T_\infty(n) = T_\infty(n/2) + \Th(\log^2 n).$$
Unrolling over $\log n$ levels, the per-level cost grows like $\log^2$ of the
current size, and the sum is dominated by the top levels; it solves to
$T_\infty(n) = \Th(\log^3 n)$.

\textbf{Parallelism.} $T_1/T_\infty = \Th(n\log n)/\Th(\log^3 n) =
\Th(n/\log^2 n)$. Compare this with the naive version's $\Th(\log n)$. We went from
``about $30$ useful processors for a billion elements'' to ``about a million''
(for $n=10^9$, $\log_2 n \approx 30$, so $n/\log^2 n \approx 10^9/900 \approx
10^6$). \emph{Same work, span slashed from $\Th(n)$ to $\Th(\log^3 n)$.}

\begin{complexitybox}
\textbf{Parallel merge sort, $n$ elements.}
\begin{center}
\begin{tabular}{lccc}
\toprule
Version & Work $T_1$ & Span $T_\infty$ & Parallelism \\
\midrule
Serial merge & $\Th(n\log n)$ & $\Th(n)$ & $\Th(\log n)$ \\
Parallel merge & $\Th(n\log n)$ & $\Th(\log^3 n)$ & $\Th(n/\log^2 n)$ \\
\bottomrule
\end{tabular}
\end{center}
The single change --- merging in parallel instead of serially --- raises the
parallelism by roughly a factor of $n/\log^3 n$. It is the chapter's signature
moral: \emph{a serial combine step can be the entire bottleneck, and parallelizing
it is where the real wins hide.}
\end{complexitybox}

\begin{focusbox}
The pattern you just watched, twice (matrix multiply, then merge sort):
\begin{enumerate}
\item Parallelize the obvious independent parts (the recursive calls). Work stays
the same; span improves only if the \emph{combine} step is cheap.
\item Discover that a leftover \emph{serial} step (the dot product; the merge)
dominates the span and chokes the parallelism.
\item Parallelize \emph{that} step too, with its own divide-and-conquer. Span
collapses; parallelism soars; work is unchanged.
\end{enumerate}
Whenever a parallel algorithm disappoints, look for the serial step hiding in the
combine. It is almost always there.
\end{focusbox}

\section{Brent's theorem, and the right way to think about it}

The last result ties the model together and tells you, before you ever run
anything, roughly how fast your algorithm will go on any machine you might buy. We
already proved the working form (Theorem~\ref{ch35:thm-greedy}); now we name it
and extract its meaning.

\begin{theorem}[Brent's theorem]\label{ch35:thm-brent}
Any computation with work $T_1$ and span $T_\infty$ can be executed on $P$
processors in time
$$T_P \le \frac{T_1 - T_\infty}{P} + T_\infty.$$
\end{theorem}

This is a slightly sharper cousin of the greedy bound (the $T_\infty$ work on the
critical path is counted once, on the span side, rather than being divided again
by $P$). For our purposes the looser $T_P \le T_1/P + T_\infty$ says everything
important, and we will reason with it.

\begin{intuitionbox}
Brent's theorem is a promise with two terms, and they pull against each other:
\begin{itemize}
\item The term $T_1/P$ is the \emph{work term}: if you had no dependencies at all,
$P$ processors would split the $T_1$ work perfectly and finish in $T_1/P$.
\item The term $T_\infty$ is the \emph{span term}: the unavoidable serial spine
that no number of processors can shorten.
\end{itemize}
For small $P$, the work term $T_1/P$ is large and dominates --- you are
\emph{work-bound}, and buying more processors helps a lot (speedup is nearly
linear). As $P$ grows, $T_1/P$ shrinks until it sinks below $T_\infty$; past that
point you are \emph{span-bound}, and buying more processors does almost nothing.
The crossover happens right around $P \approx T_1/T_\infty$ --- which is exactly
the \emph{parallelism}. That is the precise sense in which parallelism is ``the
number of processors worth buying.'' Figure~\ref{ch35:brent-curve} draws the two
regimes.
\end{intuitionbox}

\begin{figure}[h]
\centering
\begin{tikzpicture}[>=Stealth]
\draw[->,thick] (0,0) -- (7.4,0) node[right,font=\small] {$P$ (processors)};
\draw[->,thick] (0,0) -- (0,4.3) node[above,font=\small] {$T_P$ (time)};
% work term T1/P decay
\draw[sourceblue,very thick,domain=0.55:7,samples=120,smooth]
  plot (\x, {2.2/\x + 0.55});
% span floor
\draw[trapred,thick,dashed] (0,0.55) -- (7,0.55)
  node[right,font=\scriptsize,trapred] {$T_\infty$ floor};
% crossover marker
\draw[gray,dotted] (3,0) -- (3,0.55);
\node[font=\scriptsize,gray] at (3,-0.3) {$P \approx T_1/T_\infty$};
\node[sourceblue,font=\scriptsize] at (1.7,2.7) {work-bound};
\node[trapred,font=\scriptsize] at (5.6,1.0) {span-bound};
\end{tikzpicture}
\caption{Running time $T_P$ versus processor count $P$ under Brent's bound. While
the work term $T_1/P$ dominates (left), time falls like $1/P$ --- nearly linear
speedup. Once it sinks to the span floor $T_\infty$ (right), more processors stop
helping. The knee sits at the parallelism, $P \approx T_1/T_\infty$.}
\label{ch35:brent-curve}
\end{figure}

\begin{aestheticbox}
This single picture is why every chip vendor and cloud provider cares about your
algorithm's span, not just its work. Two algorithms with identical work
$\Th(n\log n)$ can have wildly different spans ($\Th(n)$ versus $\Th(\log^3 n)$ for
our two merge sorts), and that difference \emph{is} the difference between a
program that stops scaling at $30$ cores and one that keeps scaling to millions.
When someone says ``this code does not scale,'' they almost always mean its span
is too large relative to its work. You now have the vocabulary to diagnose and fix
exactly that.
\end{aestheticbox}

\begin{remark}[A reminder about what these bounds hide]
The model deliberately ignores real costs that matter in practice: the time to
\emph{communicate} data between cores, the contention for shared memory and
caches, and the overhead of the scheduler itself. A computation that is wonderful
in work and span can still be slow if every strand fights over the same cache
line, or if the spawned pieces are so tiny that scheduling them costs more than
running them (which is why real implementations stop recursing and switch to
serial code below some grain size). The work--span model tells you about an
algorithm's \emph{inherent} parallelism --- a necessary condition for scaling, and
the right first thing to analyze --- but it is the floor of the discipline, not
its ceiling.
\end{remark}

\section{The whole chapter in one breath}

You learned to describe computations with three keywords (\textbf{spawn},
\textbf{sync}, \textbf{parallel for}) and to see the result as a DAG of strands.
You measure that DAG with two numbers: \textbf{work} $T_1$ (total area, the serial
time) and \textbf{span} $T_\infty$ (critical-path depth, the infinite-processor
time). Their ratio is the \textbf{parallelism}, the most processors you can ever
use. The work and span laws give lower bounds on $T_P$; the greedy/Brent bound
$T_P \le T_1/P + T_\infty$ gives a matching upper bound, so a simple greedy
scheduler is always within a factor of two of optimal and achieves near-perfect
speedup when there is slack. Parallel loops cost a $\log$ in span for their
splitting tree. Races are the new hazard, defeated by giving each parallel strand
its own memory to write. And across matrix multiply and merge sort you saw the one
big design move again and again: \emph{parallelize the obvious parts, find the
serial combine step that dominates the span, and parallelize that too.} That is
how a $\Th(n)$ span becomes $\Th(\log^3 n)$, and how an algorithm goes from barely
parallel to embarrassingly so.

\section{Exercises}

\begin{exercise}[Warmup: read the two numbers]
For each, state the work, span, and parallelism in $\Th$ notation.
(a) A single \textbf{parallel for} of $n$ iterations, each $\Th(1)$.
(b) A serial loop of $n$ iterations, each $\Th(1)$.
(c) A balanced binary tree of $n$ leaves where every node does $\Th(1)$ work and a
node's two children are spawned in parallel.
\end{exercise}

\begin{exercise}[Span by inspection]
Consider \textsc{P-Fib}$(5)$. Without computing Fibonacci, draw the dependency
DAG, mark a longest path from the root to a leaf and back, and confirm its length
is $\Th(n)$ with $n=5$. How many total calls (strands) are there? Relate your two
counts to $T_1$ and $T_\infty$.
\end{exercise}

\begin{exercise}[The forgotten logarithm]
A student claims a \textbf{parallel for} of $n$ constant-time iterations has span
$\Th(1)$ ``because all iterations run at once.'' Explain precisely why the true
span is $\Th(\log n)$, referring to the splitting tree. Then give the span of two
\emph{nested} parallel fors, each of length $n$, with a $\Th(1)$ body.
\end{exercise}

\begin{exercise}[Spot the race]
You want the maximum of an array $A[1\ldots n]$. A colleague writes:
\textbf{parallel for} $i$: \textbf{if} $A[i] > m$ \textbf{then} $m \gets A[i]$,
where $m$ starts at $-\infty$. Exhibit a concrete input and interleaving on two
processors that returns the wrong answer. Then describe a race-free parallel
algorithm for the maximum, and give its work and span.
\end{exercise}

\begin{exercise}[Greedy bound arithmetic]
A computation has $T_1 = 10^9$ and $T_\infty = 10^4$. (a) What is its parallelism?
(b) Use $T_P \le T_1/P + T_\infty$ to estimate $T_P$ for $P = 100$, $P = 10^4$, and
$P = 10^6$. (c) For which of these $P$ is the speedup close to linear, and why does
it stop being so?
\end{exercise}

\begin{exercise}[Why the inner loop stays serial]
In \textsc{P-Matrix-Multiply}, explain what goes wrong if you make the inner
$k$-loop a \textbf{parallel for}. Then propose a \emph{race-free} way to
parallelize the dot product anyway (hint: a parallel reduction / tree of partial
sums) and state the span it would give the inner loop.
\end{exercise}

\begin{exercise}[Block multiply recurrences]
For the divide-and-conquer matrix multiply, derive $T_1(n) = 8T_1(n/2)+\Th(n^2)$
and $T_\infty(n) = T_\infty(n/2) + \Th(\log n)$ from the pseudocode structure,
explaining where each term comes from. Solve both (master theorem for the work,
unrolling for the span) and state the parallelism.
\end{exercise}

\begin{exercise}[The merge is the bottleneck]
For naive \textsc{P-Merge-Sort} (serial merge), solve $T_\infty(n) =
T_\infty(n/2) + \Th(n)$ carefully, showing the geometric sum
$n + n/2 + n/4 + \cdots = \Th(n)$. Then explain in one sentence why this single
$\Th(n)$ term, and not the recursion, is what caps the parallelism at $\Th(\log
n)$.
\end{exercise}

\begin{exercise}[Why pivot on the larger array]
In \textsc{P-Merge}, we always take the median of the \emph{longer} of the two
input ranges. Prove that with this choice each recursive call receives at most
$3/4$ of the total elements. Then explain why pivoting on the \emph{shorter}
array could fail to give a good span.
\end{exercise}

\begin{exercise}[Challenge: parallel prefix sums]
The \emph{prefix sums} (scan) of $A[1\ldots n]$ is the array $S$ with
$S[i] = A[1]+A[2]+\cdots+A[i]$. The obvious serial loop has span $\Th(n)$. Design a
parallel algorithm with work $\Th(n)$ and span $\Th(\log n)$. (Hint: an up-sweep
that builds partial sums in a tree, then a down-sweep that pushes prefixes back
down. Describe both passes and give a recurrence for each measure.)
\end{exercise}

\begin{exercise}[Challenge: when parallelism cannot help]
Give an example of a natural computation whose span equals its work, $T_\infty =
\Th(T_1)$, so its parallelism is $\Th(1)$ and no number of processors speeds it up.
Explain what feature of the problem forces the long dependency chain. (Think about
a computation where each step truly needs the previous step's result.)
\end{exercise}

\begin{exercise}[Challenge: grain size]
Real implementations stop spawning and switch to serial code once subproblems are
smaller than a \emph{grain size} $g$. Argue informally that coarsening a parallel
loop to chunks of size $g$ changes its span to $\Th(g + \log(n/g))$ while keeping
work $\Th(n)$. For what range of $g$ is the parallelism still $\Th(n/\log n)$ up to
constants, and why might a real system deliberately choose a larger $g$ anyway?
\end{exercise}
